\section{Objetivos}

\begin{frame}{Objetivos}
    \centering
    \vspace{-0.05cm}
    \begin{tikzpicture}[
        objetivo/.style={
            rectangle,
            rounded corners=6pt,
            fill=urjcred!8,
            draw=urjcred!45,
            text width=11.8cm,
            minimum height=0.92cm,
            align=center,
            font=\small\bfseries,
            inner xsep=8pt
        },
        card/.style={
            rectangle,
            rounded corners=7pt,
            minimum width=3.35cm,
            minimum height=2.45cm,
            align=center,
            font=\small,
            draw=black!25,
            inner sep=7pt
        },
        tiny/.style={
            rectangle,
            rounded corners=4pt,
            minimum width=1.62cm,
            minimum height=0.42cm,
            align=center,
            font=\scriptsize\bfseries,
            draw=black!25,
            inner sep=1.5pt
        },
        flecha/.style={-{Latex[length=2.2mm]}, thick, draw=black!45}
    ]
        \node[objetivo] (obj) at (0,2.35) {
            Desarrollar un modelo interpretable para clasificar estados emocionales\\
            a partir de señales EEG.
        };

        \node[card, fill=blue!8] (eeg) at (-4.2,0.1) {
            \includegraphics[width=2.85cm]{figs/eeg.png}\\[-1mm]
            {\bfseries Señal EEG}
        };

        \node[card, fill=green!10] (modelo) at (0,0.1) {
            {\Large\bfseries Modelo}\\[0.6mm]
            {\Large\bfseries explicable}\\[2mm]
            {\scriptsize entrenamiento + interpretación}
        };

        \node[card, fill=orange!10] (emociones) at (4.2,0.1) {
            {\bfseries Emociones}\\[2mm]
            \begin{tikzpicture}
                \node[tiny, fill=red!16, draw=red!35] at (0,0.68) {JOY};
                \node[tiny, fill=green!16, draw=green!35!black] at (0,0) {NEUTRO};
                \node[tiny, fill=blue!14, draw=blue!35] at (0,-0.68) {SAD};
            \end{tikzpicture}
        };

        \draw[flecha] (eeg) -- (modelo);
        \draw[flecha] (modelo) -- (emociones);

        \node[
            rectangle,
            rounded corners=5pt,
            fill=black!4,
            draw=black!18,
            text width=9.8cm,
            minimum height=0.58cm,
            align=center,
            font=\scriptsize
        ] at (0,-2.0) {Flujo completo: registro, características, clasificación y explicación.};
    \end{tikzpicture}
\end{frame}
